\documentclass[../../../include/open-logic-section]{subfiles}

\begin{document}

\olfileid{sth}{spine}{zf}
\olsection{$\Z$ lan $\ZF$: Tonggak Wigati}

Kanthi Landhesan, kita tekan tonggak wigati
liyane. Kita wis ngrembug teori $\Zminus$ lan
$\ZFminus$, kang kita sebut teori-teori
``dikurangi'' sawijining bab. Bab kang dikurangi
iku Landhesan. Dadi:

\begin{defn}
Teori $\Z$ nambahake Landhesan marang
$\Zminus$. Dadi aksioma-aksiomane yaiku
Ekstensionalitas, Gabungan, Pasangan, Himpunan
Pangkat, Tanpa Wates, Landhesan, lan kabeh
conto saka skema Pamisahan.

Teori $\ZF$ nambahake Landhesan marang
$\ZFminus$. Kanthi tembung liya, $\ZF$
nambahake kabeh conto saka Panggantian marang~$\Z$.
\end{defn}

Nanging, bisa wae ana pitakonan: yen Regularitas
ekivalen karo Landhesan ing $\ZFminus$, lan
dhasar kanggo Regularitas wis cetha, apa sebabe
kita kudu muter dalan lan milih Landhesan minangka
aksioma dhasar tinimbang Regularitas?

Yen alasan sejarah (bab sapa kang ngrumusake apa
lan kapan) disisihake, alasan pokoke yaiku
Landhesan bisa dirumusake tanpa nggunakake
definisi~$V_\alpha$. Definisi iku gumantung
marang kabeh asil ing \olref[recursion]{sec}:
kita kudu mbuktekake Rekursi Transfinit supaya
definisi kasebut nduweni dhasar. Nanging,
bukti Rekursi Transfinit nggunakake
\emph{Panggantian}. Mula, sanadyan Landhesan lan
Regularitas ekivalen relatif marang $\ZFminus$,
loro-lorone ora ekivalen relatif marang
$\Zminus$.

Malah, kahanane luwih abot tinimbang kang
katon saka cathetan cekak iki. Sanadyan ngluwihi
cakupan buku iki, nyatane $\Zminus$ lan $\Z$
loro-lorone banget ringkih kanggo netepake
$V_\alpha$. Mula, yen mung makarya ing $\Z$,
Regularitas (miturut rumusan kita) malah ora
nduweni \emph{makna}. Mula aksioma resmi kang
kita pilih yaiku Landhesan, dudu Regularitas.

Wiwit saiki, kita bakal makarya ing $\ZF$
(kajaba yen disebutake kanthi cetha liyane),
tanpa perlu nerangake maneh.

\end{document}
